\section{Related Work}
\label{sec:related}

\paragraph{Query-level routing and cascades.}
An early line of work treats model selection as an inference-time decision over
quality and cost for a whole query. FrugalGPT learns a cascade over heterogeneous APIs
to reduce cost \citep{chen2024frugalgpt}. AutoMix queries a smaller model
first, estimates reliability via self-verification, and escalates when needed
\citep{aggarwal2023automix}. Hybrid LLM routes between a local and a cloud
model based on predicted difficulty \citep{ding2024hybrid}. RouteLLM trains
routers from preference data and studies cross-pair transfer
\citep{ong2024routellm}. A recent survey formalizes the shared objective as
selecting a model subject to a cost budget \citep{varangotreille2025doing}.
These methods effectively abstract model selection on complete queries, but
they leave the intermediate calls inside multi-turn agent trajectories
under-explored, even though tool outputs and partial state already shape the
prefix.

\paragraph{Router benchmarks.}
RouterBench collects over 405k inference outcomes across diverse single-prompt
tasks, providing the first large-scale testbed for comparing routing algorithms
\citep{hu2024routerbench}. LLMRouterBench scales to 400k+ instances,
21 datasets, and 33 models under unified performance and cost metrics
\citep{li2026llmrouterbench}; although it includes 500 SWE-Bench issues,
SWE-Bench is configured as a single-query $0$-shot \textsc{Pass@1} task with
no agent scaffold and no intermediate routable call. RouterArena adds an open
comparison platform with difficulty levels and automated leaderboards
\citep{lu2025routerarena}. Triage makes one tier decision per SWE-bench Lite
issue using code-health features \citep{madeyski2026triage}. These benchmarks are valuable for comparing query-level routers, but even
when multi-step tasks such as SWE-Bench are included, routing reduces to
picking one model per issue from a pre-computed outcome matrix; the
intermediate retrieval, analysis, and debugging calls within an issue cannot
be routed individually to the cheapest sufficient model at each step.

\paragraph{Step-level routing.}
TRIM studies stepwise routing for multi-step math reasoning by identifying
critical reasoning steps and escalating them to larger models, and shows
that selective per-step upgrading can preserve accuracy while reducing
average cost \citep{kapoor2026trim}. Math reasoning, however, is a
relatively self-contained setting that lacks tool outputs, retrieval
contexts, shell traces, and code state, all of which are central to real
agent deployments. TRIM also derives labels from hypothetical reasoning
traces without executing the downgraded trajectory end-to-end to verify
task success. For executable long-context agent workloads the core
labeling problem persists: given the current prefix, which model tier is
cheap enough yet still preserves final task resolution?
